\chapter{Implementation}
\label{chap:implementation}

\begin{introduction}
This chapter walks through the concrete implementation decisions made throughout development to realize the architecture presented in Chapter~\ref{chap:architecture}, using them as case studies of the architecture's generalization and decoupling in practice. Each section takes one problem encountered during development, describes how it was solved, and relates the solution back to the seam it belongs to. Where a decision was later revised, that revision is described as well, since the revisions are themselves evidence of where the boundaries did and did not hold.
\end{introduction}

\section{Overview}
\label{sec:impl-overview}

The viewer did not start out decoupled. Its earliest form assumed a single backend, read its credentials from a checked-in configuration file, spawned a fixed set of actor types with type-specific code in the spawn path, and carried fire-simulation-specific logic in the generic entity code. The structure described in Chapter~\ref{chap:architecture} -- the service layer, the registration table, the two extension points, the coordinate seam -- is the result of a sequence of refactorings, each triggered by a concrete need: a second way of supplying credentials, a new entity type that did not fit the existing pattern, a rendering feature that only one entity type needed.

The sections below follow that sequence. Section~\ref{sec:impl-backend} covers moving backend specifics out of code. Section~\ref{sec:impl-streaming} covers streaming entities by geographic area and the reliability work that entailed. Section~\ref{sec:impl-worked-example} works through adding a new entity type end to end. Section~\ref{sec:impl-camera} covers the camera and interaction layer and why it was unaffected by the entity-side changes. Section~\ref{sec:impl-special} covers the one place genuinely domain-specific behaviour lives, and how it was kept behind the same seam as every other type. Section~\ref{sec:impl-discussion} reflects on which decisions generalised well and which introduced complexity.

\section{Decoupling the Data Backend}
\label{sec:impl-backend}

Three related changes moved knowledge of the backend out of the application's code and into configuration and runtime state.

\subsection{From a Configuration File to a Data Asset}
\label{sec:impl-dataasset}

Credentials and connection settings were originally read from a \texttt{Secrets.ini} file, git-ignored, parsed directly wherever they were needed. This had two problems: the file was parsed independently by more than one service, so their interpretations could drift, and configuration was expressed as text to be parsed rather than as structured data.

The first step replaced the file with a \texttt{UDittoSecretsAsset} data asset -- a first-class engine object, edited in the editor, loaded exactly once in \texttt{UDittoService::Initialize} from a fixed content path. This made connection settings structured data with a single load point. It is worth being precise about what this did and did not achieve: an early commit described the asset as ``encrypted'', but cooking a data asset into a package does not by itself encrypt it; real protection requires enabling package encryption and AES keys in the project's cryptography settings. The honest description is ``configuration as data, loaded once'', not ``encrypted secrets''.

\subsection{Runtime Credentials and the Startup Sequence}
\label{sec:impl-auth}

The second step moved credentials out of design-time configuration entirely. In the current design, the data asset holds only non-secret defaults (host, protocol, OAuth client id, the stream-start message) for developer convenience. Real credentials arrive at runtime: the user enters them on a dedicated login screen, and \texttt{UDittoService::Login} performs an OAuth2 password grant against a Keycloak endpoint. \texttt{Initialize} no longer authenticates at all.

The consequence is that authentication is asynchronous relative to the rest of startup, and several components need the result. This is coordinated entirely through delegates and a request queue, shown in Figure~\ref{fig:auth-sequence}: any REST call issued before a token lands is queued inside \texttt{UDittoService} and flushed automatically once it does, and \texttt{UWSService} does not open its socket until \texttt{UDittoService} broadcasts that an authorisation header is available. No caller has to reason about login timing; it either gets a live result or a result that arrives slightly later.

\begin{figure}[t]
  \centering
  \begin{tikzpicture}[node distance=6mm, every node/.style={font=\small}]
    \node[block, text width=26mm] (login) {login screen};
    \node[block, right=14mm of login, text width=30mm] (ditto) {\texttt{UDittoService}};
    \node[block, right=14mm of ditto, text width=24mm] (kc) {Keycloak};
    \node[block, below=12mm of ditto, text width=30mm] (ws) {\texttt{UWSService}};
    \node[block, below=8mm of ws, text width=30mm] (daemon) {\texttt{UEntityUpdateDaemon}};

    \draw[flow] (login) -- node[above,font=\footnotesize]{Login()} (ditto);
    \draw[flow] (ditto) -- node[above,font=\footnotesize]{password grant} (kc);
    \draw[flow] (kc) -- node[below,font=\footnotesize]{access + refresh token} (ditto);
    \draw[flow] (ditto) -- node[right,font=\footnotesize]{OnAuthHeaderReady} (ws);
    \draw[flow] (ws) -- node[right,font=\footnotesize]{connect + START-SEND-EVENTS} (daemon);
    \node[right=1mm of ditto.north east, font=\footnotesize\itshape, anchor=south east, yshift=6mm] {queued REST calls flushed on token};
  \end{tikzpicture}
  \caption[The authentication and startup sequence]{Authentication is asynchronous relative to startup. Components subscribe to \texttt{UDittoService}'s ``header ready'' broadcast; REST calls made before the token arrives are queued and replayed, so no caller handles login timing explicitly.}
  \label{fig:auth-sequence}
\end{figure}

Persistence of credentials is handled by a separate subsystem, \texttt{UCredentialStoreService}, backed by platform secure storage (Windows DPAPI or macOS Keychain Services). Only the host, username, and password are persisted; the short-lived OAuth access and refresh tokens are deliberately never written to disk. Instead, a fresh password grant is performed on every launch. This is a deliberate trade: it avoids the failure mode where a persisted refresh token silently goes stale between sessions, at the cost of one extra round-trip at startup, and it means a stored login stays usable until the user explicitly logs out. On a platform with no secure-storage backend implemented, loading simply fails and the login screen is shown every launch -- a graceful degradation rather than an error.

The login flow is isolated at the level design as well: it runs under its own minimal game mode with no pawn, no world geometry, and no Cesium globe, so none of the tile-streaming machinery starts before the user is authenticated.

\subsection{A Single Source of Truth for Connection Settings}

With settings loaded once by \texttt{UDittoService}, the other subsystems that need them read them from its getters rather than loading the asset a second time. \texttt{UWSService} in particular takes its host, protocol, and stream-start message this way. The result is that there is one interpretation of the connection configuration in the running application, and one moment at which the live connection is considered ready, rather than each subsystem arriving at its own.

\subsection{What Backend Replacement Would Still Require}
\label{sec:impl-replace-caveat}

It is important to be precise about the sense in which the backend is ``replaceable''. What has been achieved is \emph{confinement}: no code above the service layer names an HTTP verb, a URL, an OAuth scope, a WebSocket frame, or Ditto's envelope format, so the blast radius of a backend change is one layer. What has \emph{not} been done is to factor \texttt{UDittoService} behind an abstract interface, or to build a second backend implementation or an offline mock. A genuine swap today would mean rewriting the internals of \texttt{UDittoService}, \texttt{UWSService}, and the envelope parser in \texttt{UEntityUpdateDaemon} against the new backend, while keeping their public signatures. That is a bounded, well-localised task -- which is the point of the seam -- but it is not yet a drop-in replacement. This is revisited in Chapter~\ref{chap:conclusions}.

\section{Entity Streaming and Reliability}
\label{sec:impl-streaming}

\subsection{Streaming Entities by Geographic Area}

Loading every entity in the dataset regardless of where the camera is looking does not scale, and baking in an assumption about dataset size would undermine the goal of generalising across infrastructure of arbitrary extent. Instead, \texttt{ADT4MOBGamemode} maps the camera's altitude to a discrete zoom level and, above a threshold zoom, fetches only the entities within the 3$\times$3 grid of geographic tiles around the camera, expressed as quadkey ranges and queried through \texttt{UDittoService::GetThingsByGeotileBounds}. Below the threshold zoom, tile filtering is disabled and everything loads, since at that range the whole area of interest is on screen anyway. Figure~\ref{fig:tile-flow} shows the per-check decision path.

\begin{figure}[t]
  \centering
  \begin{tikzpicture}[node distance=5mm, every node/.style={font=\small}]
    \node[block, text width=40mm] (raw) {camera $\rightarrow$ lat/lon/alt\\{\footnotesize raw zoom = f(altitude)}};
    \node[block, below=of raw, text width=52mm] (hyst) {zoom-change hysteresis\\{\footnotesize new zoom must hold ${\sim}2$\,s before it commits}};
    \node[block, below=of hyst, text width=52mm] (thresh) {zoom $<$ threshold?};
    \node[block, right=10mm of thresh, text width=34mm] (all) {drop tile filter,\\keep vehicle-exclude filter};
    \node[block, below=of thresh, text width=52mm] (same) {same 3$\times$3 tile set as loaded?};
    \node[block, below=of same, text width=52mm] (deb) {debounce: camera stable ${\sim}0.75$\,s};
    \node[block, below=of deb, text width=52mm] (refresh) {refresh only changed tiles\\{\footnotesize fetch new @ each geotile precision; unload gone; protect in-view/open as orphans}};
    \draw[flow] (raw) -- (hyst);
    \draw[flow] (hyst) -- (thresh);
    \draw[flow] (thresh) -- node[above,font=\footnotesize]{yes} (all);
    \draw[flow] (thresh) -- node[right,font=\footnotesize]{no} (same);
    \draw[flow] (same) -- node[right,font=\footnotesize]{no} (deb);
    \draw[flow] (deb) -- (refresh);
  \end{tikzpicture}
  \caption[The tile-refresh decision path]{The per-check tile-refresh decision path in \texttt{ADT4MOBGamemode}. Each guard -- hysteresis, threshold, no-op detection, debounce, changed-tiles-only refresh, orphan protection -- was added in response to an observed misbehaviour.}
  \label{fig:tile-flow}
\end{figure}

Making this stable in the face of a continuously moving camera required several guards, each added in response to an observed misbehaviour:

\begin{itemize}
    \item a periodic check (every ${\sim}0.5$\,s, rather than per-frame) detects when the camera has entered a new tile;
    \item a debounce delays the actual refresh until the camera has held still for ${\sim}0.75$\,s, so sweeping the view across several tiles issues one refresh, not one per tile crossed;
    \item a refresh only destroys and re-fetches the tiles that actually changed, keeping the ones still in view;
    \item a zoom-change hysteresis (Listing~\ref{lst:hysteresis}) requires a new zoom level to hold for ${\sim}2$\,s before it is committed, because camera-altitude jitter sitting exactly on a zoom boundary would otherwise oscillate the level, and every flip rebuilds the entire scene;
    \item entities that are still in view or have an open detail window when their tile unloads are kept alive as ``orphans'' and swept later once they are neither, rather than being destroyed out from under the user.
\end{itemize}

\begin{lstlisting}[style=cpp,float=tp,caption={[The zoom-change hysteresis guard]The zoom-change hysteresis in \texttt{ADT4MOBGamemode::CheckAndRefreshTiles}. \texttt{AltitudeToZoomLevel} rounds a logarithm, so an altitude on a zoom boundary flips on sub-metre jitter; each flip would run \texttt{DestroyAllActors}.},label={lst:hysteresis}]
int32 Zoom = RawZoom;
if (LoadedZoom >= 0 && RawZoom != LoadedZoom)
{
    if (RawZoom == CandidateZoom)
    {
        ZoomHoldTimer -= DeltaSeconds;
        if (ZoomHoldTimer > 0.f)
            Zoom = LoadedZoom;          // not held long enough yet -- stay put
    }
    else
    {
        CandidateZoom = RawZoom;
        ZoomHoldTimer = ZoomChangeHoldSeconds;
        Zoom = LoadedZoom;
    }
}
else { CandidateZoom = -1; }
\end{lstlisting}

Different data sources in the backend were also found to encode an entity's geotile at different precisions -- simulated vehicles at one zoom, imported static infrastructure at another -- so every tile query is issued at each known precision and the results merged, rather than assuming one encoding.

The live event subscription is scoped to match: when zoomed out past the tile-filtering threshold, individual simulated vehicles are sub-pixel and their high-churn update stream would dominate the WebSocket load for no visible benefit, so the subscription filter excludes that namespace entirely at that range.

\subsection{Serialising Concurrent Tile Requests}
\label{sec:impl-serialise}

An intermittent failure appeared during this work: \texttt{GetThingsByGeotileBounds} would occasionally log a transport-level failure -- not an HTTP error status, but the request failing to complete at all. It was traced to the refresh logic firing every changed tile's request concurrently in a tight loop, which the backend host handled unreliably under bursts. The fix at the time was to serialise the fetches: a queue held the pending tiles and started each tile's request only from the previous one's completion callback.

This was later removed. The residual flakiness was confirmed to be a server-side problem with the specific backend host, present before this code existed and independent of it, and once the backend was fixed the burst behaviour was no longer an issue and the current code fires the per-tile requests directly again. The episode is a useful case study for Section~\ref{sec:impl-discussion}: a decoupled architecture still has to respect the operational limits of whatever backend sits behind it, and a workaround written against one backend's transient behaviour is itself a form of coupling -- one that had to be carried until the backend changed and then deliberately unwound.

\subsection{Per-Entity Expiry for Simulated Fleets}
\label{sec:impl-ttl}

Ephemeral entities carry an expiry timestamp. For simulated vehicle fleets, the data source stamps the entire fleet with one shared timestamp, so honouring it literally destroys every vehicle in the same frame when the deadline passes, regardless of which ones are still being simulated. The generic path could not express ``expire these individually''. This is handled through the stateless extension point: the vehicle type's \texttt{UEntityTypeExtension} returns a per-entity staleness timeout (60\,s) instead, so an actor is removed only when no update for that specific entity has arrived within that window, with the timer reset on every message. Entities that genuinely stopped updating are removed, staggered naturally in time; entities still being simulated are not. The same extension also opts the type out of the per-message reflected-struct rebuild (\texttt{NeedsStructSyncOnLiveUpdate}), since a simulated vehicle's position is read straight from its raw JSON and the reflection walk -- a full traversal of the struct per message -- would be pure overhead at that update frequency. Both decisions are visible in Listing~\ref{lst:vehicleext} in the previous chapter.

\section{Adding an Entity Type: A Worked Example}
\label{sec:impl-worked-example}

To make the generalisation concrete, consider adding a hypothetical new infrastructure type -- say, an air-quality sensor with a \texttt{thingId} containing \texttt{air-quality} and attributes for the current readings. The full set of changes is:

\begin{enumerate}
    \item \textbf{Schema.} Add \texttt{AirQualityStruct.h} with \texttt{USTRUCT}s mirroring the sensor's JSON: a top-level struct with the latitude, longitude, and an \texttt{attributes} sub-struct holding the readings. \texttt{FJsonObjectConverter} handles deserialisation from that declaration alone.
    \item \textbf{Registration.} Add one line to \texttt{RegisterAllEntityTypes}: \texttt{Factory.RegisterType("air-quality", FAirQualityData::StaticStruct(), TEXT("Air Quality Sensor"), \ldots)}.
    \item \textbf{(Optional) presentation.} If the default grey badge and empty info panel are not wanted, add a \texttt{UAirQualityExtension} under \texttt{EntityDependencies/AirQuality/} overriding \texttt{GetBadgeColor}, \texttt{GetBadgeLabel}, and \texttt{GetDefaultInfoFields}.
\end{enumerate}

Nothing else changes. The factory now spawns an \texttt{ATempUIActor} for any matching \texttt{thingId}; the tile-streaming loop fetches it like any other entity; \texttt{UEntityUpdateDaemon} routes its live updates; the actor patches its JSON, updates its position through the coordinate seam, and refreshes the UI; selection, hover, and both camera modes work unmodified. There is no edit to the spawn pipeline, the update loop, the renderer, or the camera. This three-step (often two-step) process is the operational meaning of the generalisation claim in RQ1.

\section{Camera and Interaction}
\label{sec:impl-camera}

The camera is a single pawn, \texttt{AUnifiedPawn}, with two modes selected by an \texttt{ECameraMode} enum. In the RTS (overhead) mode it holds a fixed downward pitch with a spring arm whose length is the zoom, ranging from roughly $1.5$\,km to planetary distance, and re-aligns itself to the globe's local north and east every tick via \texttt{UpdateRtsAlignment}. In the free-fly mode it collapses the spring arm for full six-degree-of-freedom movement, staying upright relative to the globe surface by querying the local up-vector from the georeference in \texttt{UpdateFreeFlyAlignment}. Both modes' globe alignment is computed by projecting the ECEF north direction onto the local tangent plane.

The point worth making about the camera is what it does \textit{not} reference. It queries the globe -- the georeference -- for orientation, and it operates on actors and on the cursor for interaction. It has no knowledge of entity types, of the backend, or of the data model at all. Selection and hover are likewise generic: \texttt{USelectionManager}, a local-player subsystem, broadcasts hover and selection changes, and \texttt{ATempUIActor} responds by toggling custom depth-stencil values on all of its primitive components, with no per-type branching; the HUD widget binds to the same broadcasts to show or hide the entity detail window. This is precisely why the entire camera and interaction layer needed no changes while the entity side was being repeatedly restructured -- it was never coupled to the entity side in the first place.

\section{Special-Case Handling}
\label{sec:impl-special}

Fire-simulation entities are the one place genuinely infrastructure-specific behaviour exists in the viewer. A fire entity loads two separate multi-mesh models (a spread cone and a simulation body), fetches and parses GeoJSON perimeter geometry, and projects a precise terrain-exclusion polygon onto the globe so the burned area reads correctly against the terrain.

None of this could be expressed through a table entry or a stateless extension, because it involves live per-instance state: fetched GeoJSON, resolved model URLs, parsed vertex rings. It is therefore implemented as a \texttt{UEntityBehaviorComponent} subclass, \texttt{UFireBehaviorComponent} (Listing~\ref{lst:fire}), returned by the fire type's extension and attached generically by \texttt{ATempUIActor::Initialize} to each fire actor. The actor does not know what the component does; it just knows that some types supply one. The component overrides the two narrow points it needs -- it declares that it handles its own model loading, and it supplies exact exclusion polygons in place of the actor's default behaviour of hulling the visible mesh vertices.

\begin{lstlisting}[style=cpp,float=tp,caption={[The fire behaviour component]\texttt{UFireBehaviorComponent}. It is the only entity type that reaches the third rung of the entity model. The generic actor attaches it without knowing what it does.},label={lst:fire}]
UCLASS()
class DT4MOB_API UFireBehaviorComponent : public UEntityBehaviorComponent
{
    GENERATED_BODY()
public:
    virtual void OnEntityInitialized() override;   // load cone/sim GLBs, fetch GeoJSON
    virtual void OnEntityDataChanged() override;   // re-fetch on live perimeter update
    virtual bool HandlesOwnModelLoading() const override { return true; }
    virtual bool GetExclusionPolygons(
        TMap<FString, TArray<FVector2D>>& OutPolygons) const override;  // "Cone", "Simulation"
};
\end{lstlisting}

This section also carries a two-step history that illustrates the generalisation boundary being found in practice, sketched in Figure~\ref{fig:fire-fix}. The first refactoring separated the staged model loads for special-logic entities like fire out of the generic entity-load path, so that the generic path no longer had a fire-shaped branch in it. A later fix addressed a subtler coupling: the fire cone and the simulation exclusion polygon had been mutually exclusive on screen, not by design but as a side effect of sharing a single piece of state -- one \texttt{ACesiumCartographicPolygon} slot. The exclusion mechanism was changed so the behaviour component returns a \emph{map} of named polygons and \texttt{ATempUIActor} spawns one Cesium polygon per entry, so the cone and the simulation area can be shown at once. Both fixes pushed fire-specific behaviour further behind the same seam every other type uses, rather than special-casing it in the factory or the update loop.

\begin{figure}[t]
  \centering
  \begin{tikzpicture}[node distance=7mm, every node/.style={font=\small}]
    \node[block, text width=44mm] (b1) {\textbf{before}\\{\footnotesize one polygon slot shared}\\cone \emph{xor} simulation visible};
    \node[block, right=16mm of b1, text width=48mm] (a1) {\textbf{after}\\{\footnotesize \texttt{TMap<FString,TArray<FVector2D>>}}\\one Cesium polygon spawned per entry\\cone \emph{and} simulation visible};
    \draw[flow] (b1) -- (a1);
  \end{tikzpicture}
  \caption[The exclusion-polygon coupling fix]{Removing an accidental coupling: the cone and simulation areas were mutually exclusive only because they shared one polygon slot. Keying them by name in a map removed the constraint without any change to the generic actor beyond iterating the map.}
  \label{fig:fire-fix}
\end{figure}

\section{Discussion}
\label{sec:impl-discussion}

Some of these decisions generalised cleanly and some introduced complexity that is worth naming.

The registration table plus the two extension points held up well. Every entity type added after the pattern was in place fit one of its three rungs, and the generic spawn, update, and render paths did not need to change to accommodate any of them -- the worked example in Section~\ref{sec:impl-worked-example} is the normal case. The fire component -- the most demanding entity in the system -- fit the stateful-component rung without forcing an exception. The coordinate-conversion seam and the configuration-as-data change were similarly low-cost and paid off immediately, the latter by removing an entire class of ``which service parsed the config'' bug.

The decoupling from the backend is real but partial, as Section~\ref{sec:impl-replace-caveat} sets out: backend knowledge is confined to one layer, but that layer is a concrete implementation, not an interface with alternatives behind it, and no offline mock was built. Backend independence is therefore demonstrated by construction and by inspection -- one can point to every place the backend is named and observe that they are all in the service layer -- rather than by a working substitution.

The tile-streaming loop is where complexity accumulated. Each of its guards -- debounce, changed-tiles-only refresh, zoom hysteresis, orphan sweep, multi-precision geotile queries -- is individually justified by an observed failure, but together they make that loop the most intricate part of the codebase, and its correctness rests on empirically determined constants (hold times, thresholds) rather than on a clean invariant. The request-serialisation workaround, added and then removed, is the clearest example of the architecture being unable to fully insulate itself from backend behaviour: for a period, the viewer's fetch logic was shaped by one host's transient reliability rather than by its own needs.

This sets up the evaluation in Chapter~\ref{chap:evaluation}. The performance benchmarks target the entity spawning and rendering path -- exercised most heavily by high-cardinality types such as signs and simulated vehicles -- to establish how it scales with count and which rendering strategy such a type should use (RQ3); the per-type indirection itself (one map lookup for the extension, one for the schema, per spawn) is negligible by inspection and is not separately measured. The user testing targets the camera and interaction layer, since the question there is whether a layer built to be generic over the data still feels coherent to someone using it to explore a specific dataset (RQ4).
